\begin{table*}[t]
\centering
\caption{实验二的设计归因。效应倍率先在同一片段或 event 内计算 Component off / on，再报告 median [Q1, Q3]；所有指标越低越好，倍率大于 1 表示启用设计更优。VCD 行的 on/off 分别表示分数排序的 AURC 与忽略排序的全覆盖风险，并非冻结阈值接纳消融；时序策略行比较均关闭 StaticLock 的 Linear/SLERP 与 Smoothed KF。}
\label{tab:exp2-final}
\small
\setlength{\tabcolsep}{3.0pt}
\resizebox{\textwidth}{!}{%
\begin{tabular}{llllll}
\toprule
组件关闭 & 参照指标 & EgoAnchor (on) & Component off & 效应倍率 & 一致性 \\
\midrule
采集时刻对齐 & 同候选复合 P95 & Capture time 15.38 [13.44, 19.13]~mm & Arrival time 43.75 [38.14, 55.60]~mm & 2.69 [2.38, 3.00]$\times$ & 12/12 \\
StaticLock & 中心化静止 P95 & 0.82 [0.58, 0.92]~mm & 13.73 [11.90, 14.89]~mm & 17.06 [13.89, 27.52]$\times$ & 12/12 \\
VCD 判别性 & event AURC & 5.65 [4.90, 5.92]~mm & Full coverage 6.53 [5.98, 7.31]~mm & 1.21 [1.18, 1.24]$\times$ & 12/12 \\
时序策略 & 平移 / 旋转 aligned RMSE & \shortstack{Linear/SLERP (StaticLock off)\\8.93 [7.41, 11.61]~mm / 4.70 [4.35, 4.82]$^\circ$} & \shortstack{Smoothed KF (StaticLock off)\\45.64 [37.58, 49.73]~mm / 12.21 [11.28, 13.77]$^\circ$} & \shortstack{T: 5.46 [4.02, 5.73]$\times$\\R: 2.63 [2.44, 2.78]$\times$} & \shortstack{T: 16/16\\R: 12/12} \\
\bottomrule
\end{tabular}%
}
\end{table*}
